\documentclass{article}
\usepackage{amsmath,amssymb,amsthm}
\IfFileExists{hyperref.sty}{%
  \IfFileExists{color.sty}{\usepackage{color}\definecolor{paperlinkblue}{rgb}{0.10,0.25,0.42}%
    \usepackage[colorlinks=true,linkcolor=paperlinkblue]{hyperref}}{\usepackage[hidelinks]{hyperref}}%
}{}
\ifdefined\hyperlink
  \newcommand{\paperanchor}[1]{\hypertarget{#1}{}}%
  \newcommand{\paperlink}[2]{\hyperlink{#1}{\begin{NoHyper}#2\end{NoHyper}}}%
\else
  \ifdefined\pdfextension
    \ifnum\outputmode>0
      \def\paperdest{\pdfextension dest}\def\paperstart{\pdfextension startlink}%
      \def\paperend{\pdfextension endlink}\def\papercolor{\pdfextension colorstack}%
    \fi
  \else\ifdefined\pdfstartlink
    \ifnum\pdfoutput>0
      \let\paperdest\pdfdest\let\paperstart\pdfstartlink\let\paperend\pdfendlink\let\papercolor\pdfcolorstack
    \fi
  \fi\fi
  \ifdefined\paperdest
    \newcommand{\paperanchor}[1]{\paperdest name{#1} xyz}%
    \newcommand{\paperlink}[2]{\leavevmode\paperstart attr{/Border [0 0 0]} goto name{#1}%
      \papercolor 0 push {0.10 0.25 0.42 rg 0.10 0.25 0.42 RG}#2\papercolor 0 pop\paperend}%
  \else\ifdefined\XeTeXversion
    \newcommand{\paperanchor}[1]{\special{pdf:dest (#1) [@thispage /XYZ null @ypos null]}}%
    \newcommand{\paperlink}[2]{\leavevmode\special{pdf:bann << /Subtype /Link /Border [0 0 0] /A << /S /GoTo /D (#1) >> >>}%
      \special{color push rgb 0.10 0.25 0.42}#2\special{color pop}\special{pdf:eann}}%
  \else
    \PackageError{lean-to-latex}{No PDF hyperlink backend available}{Use XeLaTeX, pdfLaTeX, LuaLaTeX, or install hyperref.}%
  \fi\fi
\fi
\newcommand{\paperlineanchor}[1]{\raisebox{10pt}[0pt][0pt]{\paperanchor{#1}}}
\DeclareRobustCommand{\paperresultname}[1]{\paperlineanchor{theorem:\arabic{theorem}}#1}
\newtheoremstyle{paperresult}{8pt}{5pt}{\normalfont\raggedright}{}{\bfseries}{.}{\newline}{}
\theoremstyle{paperresult}
\newtheorem{theorem}{\paperresultname{Theorem}}
\newcounter{paperpremise}
\newtheorem{lemma}[theorem]{\paperresultname{Lemma}}
\newtheorem{corollary}[theorem]{\paperresultname{Corollary}}
\newcounter{definingequation}
\newenvironment{proofstep}{\par\smallskip\ifdim\leftskip<3em\relax\advance\leftskip by 1em\relax\fi\noindent\ignorespaces}{\par\smallskip}
\setlength{\textwidth}{6.5in}
\setlength{\oddsidemargin}{0in}
\AddToHook{env/proof/begin}{\raggedright}
\setlength{\emergencystretch}{\linewidth}
\title{Two arithmetic bounds}
\author{}
\date{}
\begin{document}
\maketitle
\section{Results}
\begin{theorem}[{Weighted bound}]\label{theorem:1}
Let \(x\) and \(y\) be rational numbers. Assume \(2 \cdot x + y \le 9\). Assume \(1 \le y\). Then \(x \le 4\).
\end{theorem}

\begin{proof}
Suppose for contradiction that \(x > 4\).
Subtracting the right side of \(2 \cdot x + y \le 9\) gives \(2 \cdot x + y - 9 \le 0\).
Subtracting the right side of \(1 \le y\) gives \(1 - y \le 0\).
Subtracting the right side of \(x > 4\) gives \(4 - x < 0\).
Multiplying \(4 - x < 0\) by 2 gives \(2 \cdot \left(4 - x\right) < 0\).
Adding \(2 \cdot x + y - 9 \le 0\), \(1 - y \le 0\) and \(2 \cdot \left(4 - x\right) < 0\) gives \(2 \cdot x + y - 9 + \left(1 - y\right) + 2 \cdot \left(4 - x\right) < 0\).
The terms on the left cancel, so the inequality would give \(0 < 0\), a contradiction.
\end{proof}

\begin{theorem}[{Rounded bound}]\label{theorem:2}
Let \(x\) and \(y\) be integers. Assume \(2 \cdot x \le y\). Assume \(y \le 9\). Then \(x \le 4\).
\end{theorem}

\begin{proof}
Suppose for contradiction that \(\neg \left(x \le 4\right)\).
Writing \(2 \cdot x \le y\) as an integer constraint gives \(y - 2\,x \ge 0\).
Taking the linear combination of \(y - 2\,x \ge 0\) and \(y \le 9\) with respective coefficients 1 and -1 gives \(-2\,x \ge -9\).
Divide \(-2\,x \ge -9\) by \(-2\), reversing the order. Since \(x\) is integral, \(x \le \left\lfloor\frac{9}{2}\right\rfloor = 4\).
From \(\neg \left(x \le 4\right)\) we get \(4 < x\), hence \(x \ge 5\) over the integers.
Combining the bounds gives \(5 \le x \le 4\).
This is impossible because \(5 > 4\).
\end{proof}

\end{document}
